\documentclass[11pt,a4paper]{article}

\usepackage[T1]{fontenc}
\usepackage[utf8]{inputenc}
\usepackage{lmodern}
\usepackage{microtype}
\usepackage[margin=2.2cm]{geometry}
\usepackage{booktabs}
\usepackage{array}
\usepackage{longtable}
\usepackage{enumitem}
\usepackage[hidelinks]{hyperref}

\setlist{nosep, leftmargin=*}
\setlength{\parskip}{0.5em}
\setlength{\parindent}{0pt}

\newcommand{\code}[1]{\texttt{#1}}
\newcolumntype{L}[1]{>{\raggedright\arraybackslash}p{#1}}

\begin{document}

\begin{center}
  {\LARGE \textbf{cgrad}: A Minimal Neural Network Library\\[0.2em]
  with Automatic Differentiation in C\par}
  \vspace{0.8em}
  {\large Gregory Frommelt\par}
  \vspace{0.2em}
  {\small Project Proposal \quad Programming Languages and
  Paradigms \quad University of Zurich, Fall 2026\par}
\end{center}

\section{Project Description}

\textbf{cgrad} is a small, self-contained library for building and training
feed-forward neural networks, written from scratch in C. It uses only the C
standard library and \code{libm}, with no machine-learning or linear-algebra
libraries. At its core is a tensor type and a computational graph that is built
dynamically during the forward pass. A reverse-mode automatic differentiation
engine traverses this graph backwards to compute gradients. On top of it, the
library provides common differentiable operations, simple network layers and a
gradient-based optimiser.

The library will be demonstrated end to end. First it learns small toy problems
(e.g.\ XOR). Then it trains a multi-layer perceptron on a standard real-world
benchmark (e.g.\ handwritten-digit classification on MNIST, with the dataset
loaded directly from its binary file format).

\section{Motivation: Why C?}

\begin{itemize}
  \item \textbf{Explicit lifetimes.} An autograd engine creates and discards
    many graph nodes in every training iteration. In garbage-collected
    languages this is invisible. In C, the lifetime and ownership of every node
    has to be designed explicitly, which makes the core engineering problem of
    such a system tangible.
  \item \textbf{Control over memory and performance.} Numerical code depends
    on contiguous memory, data layout and cache behaviour. C exposes all of
    these directly, and it is the layer on which established numerical and ML
    libraries are built.
  \item \textbf{Polymorphism by hand.} Each operation needs its own backward
    rule. In C this is built from structs and function pointers, which shows
    how dynamic dispatch works beneath object-oriented abstractions.
  \item \textbf{No hidden runtime.} Every allocation and every computation is
    explicit. This makes the system easy to reason about, to measure and to
    compare with high-level frameworks.
\end{itemize}

\section{Key Language Features}

The project will likely rely on the following C features:

{\small
\begin{longtable}{@{}L{0.27\textwidth} L{0.33\textwidth} L{0.34\textwidth}@{}}
\toprule
\textbf{Feature} & \textbf{Why} & \textbf{How} \\
\midrule
\endhead
Structs, \code{typedef}, \code{enum}
  & Model the core abstractions
  & Tensors, graph nodes, layers, optimiser state \\
Pointers, pointers to pointers
  & Represent graph edges and collections
  & Links from a node to its inputs; lists of nodes and parameters \\
Dynamic memory (\code{malloc}/\code{free})
  & Graphs are built and destroyed at run time
  & Explicit ownership rules for long-lived and per-iteration objects \\
Function pointers
  & Per-operation gradient rules
  & Each node stores the backward function of the operation that created it \\
\code{void *}
  & Operation-specific saved state
  & Generic context attached to graph nodes \\
Headers, \code{static}, opaque types
  & Encapsulation without classes
  & Small modules with clear public interfaces and private internals \\
Preprocessor
  & Build configuration and testing
  & Header guards, debug/release switches, assertion macros for tests \\
Binary I/O, fixed-width integers
  & Read datasets and possibly store models
  & Parsing binary file formats in a portable way \\
\bottomrule
\end{longtable}
}

\section{System Architecture}

The library will be split into a few loosely coupled modules, roughly:

\begin{itemize}
  \item \textbf{Tensor core}: a two-dimensional ($n \times m$) numeric container with its data,
    gradient storage and basic utilities.
  \item \textbf{Autodiff engine}: records the computational graph, orders it
    topologically and runs the backward pass.
  \item \textbf{Operations}: differentiable building blocks such as arithmetic,
    matrix multiplication, activation functions and loss functions, each with
    a forward and a backward rule.
  \item \textbf{Network layers}: e.g.\ fully connected layers composed into a
    multi-layer perceptron.
  \item \textbf{Optimisation}: parameter updates based on the computed
    gradients.
  \item \textbf{Data handling}: dataset loading, batching and shuffling.
\end{itemize}

\textbf{Memory strategy.} Objects fall into two lifetimes. Model parameters
persist across the whole training run. Intermediate results of the forward pass
live for a single iteration and are released together after each update step.
Clear ownership rules decide which component allocates and which one frees each
kind of object.

\textbf{Tooling.} The project uses C11 and a plain \code{Makefile} with strict
compiler warnings, a small hand-written test suite, and runtime checks for
memory errors and undefined behaviour (sanitizers and leak detection).

\section{Scope}

\textbf{Core (must-haves)}
\begin{itemize}
  \item Tensor type with a dynamically built computational graph
  \item Reverse-mode automatic differentiation
  \item A set of common differentiable operations (arithmetic, matrix
    operations, activations, losses)
  \item A multi-layer perceptron and at least one gradient-based optimiser
  \item Verification of the computed gradients (e.g.\ against numerical
    approximations)
  \item Training on a toy problem and on a real dataset
\end{itemize}

\textbf{Possible extensions (chosen depending on progress)}
\begin{itemize}
  \item Additional optimisers
  \item Custom memory management for graph nodes (e.g.\ arena or pool allocation)
  \item Performance work (e.g.\ cache-aware kernels, multi-threading)
  \item Saving and loading trained models
  \item Comparison with an established framework in accuracy and speed
\end{itemize}

\textbf{Out of scope}
\begin{itemize}
  \item GPU support
  \item Convolutional or recurrent architectures; tensors beyond two dimensions
  \item General-purpose broadcasting
  \item Bindings to other languages
\end{itemize}

\section{Challenges and Mitigation}

{\small
\begin{longtable}{@{}L{0.27\textwidth} L{0.33\textwidth} L{0.34\textwidth}@{}}
\toprule
\textbf{Challenge} & \textbf{Mitigation} & \textbf{Scope fallback} \\
\midrule
\endhead
Memory management of the graph
  & Ownership rules defined up front; sanitizers and leak checks from the start
  & Keep the simplest working allocation strategy \\
Correctness of gradients
  & Numerical gradient checks for every operation
  & Reduce the set of supported operations \\
Numerical stability
  & Numerically stable formulations of critical functions (e.g.\ softmax and
    logarithms)
  & Restrict to well-behaved operations \\
Training performance
  & Mini-batching, optimised builds, attention to memory layout
  & Train on a subset of the data or a smaller model \\
Learning curve of C
  & Incremental development, small test-covered steps
  & Drop optional extensions first \\
\bottomrule
\end{longtable}
}

Core functionality comes first. Extensions are only started once the core
works end to end, so the project stays deliverable even if parts take longer
than expected.

\end{document}
